\section*{APPENDIX A. IMPLEMENTATION AND DATA REFERENCE}
\addcontentsline{toc}{section}{\protect\numberline{}APPENDIX A. IMPLEMENTATION AND DATA REFERENCE}
\setcounter{section}{6}
\setcounter{subsection}{0}
\setcounter{figure}{0}
\setcounter{table}{0}

\subsection{Radio Simulator Module Mapping}

\begin{longtable}{p{3.2cm}p{10.7cm}}
\toprule
Module & Research responsibility\\
\midrule
\endfirsthead
\toprule
Module & Research responsibility\\
\midrule
\endhead
\texttt{main.py} & Reads scenario inputs and starts the interactive simulation.\\
\texttt{config.py} & Defines radio, mobility, timing, and output constants.\\
\texttt{hexagons.py} & Generates axial hexagonal gNB locations.\\
\texttt{mobility.py} & Implements roads, intersections, turns, boundaries, and stop-and-go motion.\\
\texttt{radio\_models.py} & Implements LOS probability, UMa/UMa-AV path loss, fading, RSRP, interference, SINR, and association.\\
\texttt{simulation.py} & Owns simulation state, executes each time step, updates visualization, and coordinates logging.\\
\texttt{logger.py} & Builds the per-step dataset and writes the combined CSV.\\
\texttt{ue.py} & Validates and stores aerial-UE metadata.\\
\texttt{utils.py} & Reads numeric input and applies defaults.\\
\bottomrule
\end{longtable}

\subsection{Protocol Test Environment}

\begin{longtable}{p{3.2cm}p{10.7cm}}
\toprule
Component & Research responsibility\\
\midrule
\endfirsthead
\toprule
Component & Research responsibility\\
\midrule
\endhead
\texttt{free5GC} & Provides the 5G Core functions used by the protocol-level experiment, especially AMF-side NGAP/N2 connectivity and user-plane support through GTP-U/N3.\\
\texttt{StormSim} & Emulates UE and gNB behavior for selected handover scenarios, connects to the 5G Core, injects configurable loss, delay, and jitter, and exports trial-summary and timer-event logs.\\
\texttt{Python radio simulator} & Generates radio-level mobility and measurement traces used to inspect RSRP, SINR, serving-cell changes, and neighbor-cell behavior.\\
\bottomrule
\end{longtable}

\subsection{StormSim Module Mapping}

\begin{longtable}{p{4.2cm}p{9.7cm}}
\toprule
Package & Research responsibility\\
\midrule
\endfirsthead
\toprule
Package & Research responsibility\\
\midrule
\endhead
\texttt{cmd/emulator} & Selects normal, replay, CSV, measurement, failure, and CHO modes.\\
\texttt{pkg/config} & Loads YAML and CSV input schemas.\\
\texttt{pkg/model} & Defines shared UE, gNB, session, event, RRC, and CHO data.\\
\texttt{internal/common/fsm} & Executes asynchronous state machines.\\
\texttt{internal/common/logger} & Records protocol, timer, delay, and CSV/NDJSON outputs.\\
\texttt{internal/transport/rlink} & Emulates UE--gNB channels with directional loss, delay, jitter, and counters.\\
\texttt{internal/transport/sctpngap} & Provides SCTP/NGAP transport to the AMF.\\
\texttt{internal/core/uecontext} & Implements UE 5GMM/5GSM, NAS security, sessions, handover execution, recovery timers, and CHO evaluation.\\
\texttt{internal/core/gnbcontext} & Implements gNB context, NGAP handling, Xn/N2 handover, path switch, and CHO preparation.\\
\texttt{internal/core/monitor} & Tracks handover phases, results, groups, NTN visibility metadata, and OAM responses.\\
\texttt{internal/scenarios} & Creates UE/gNB experiments and exports trial/timer results.\\
\texttt{monitoring} & Provides OAM, resource, PCAP, and GTP-U support.\\
\bottomrule
\end{longtable}

\subsection{Python Radio CSV Dictionary}

For UE index \(i\), the logger creates the following fields.

\begin{longtable}{p{4.6cm}p{9.3cm}}
\toprule
Column pattern & Meaning\\
\midrule
\endfirsthead
\toprule
Column pattern & Meaning\\
\midrule
\endhead
\texttt{Step} & Discrete simulation frame.\\
\texttt{ue\(i\)\_x}, \texttt{ue\(i\)\_y} & UE Cartesian position in metres.\\
\texttt{ue\(i\)\_height} & UE height above ground in metres.\\
\texttt{ue\(i\)\_direction} & 0: east, 1: north, 2: west, 3: south.\\
\texttt{ue\(i\)\_connected\_bs} & Serving gNB array index.\\
\texttt{ue\(i\)\_los\_probability} & Calculated LOS probability for the serving link.\\
\texttt{ue\(i\)\_los\_state} & Cached sampled state: LOS or NLOS.\\
\texttt{ue\(i\)\_pathloss} & Path loss before shadow fading, in dB.\\
\texttt{ue\(i\)\_shadow\_fading} & Additional shadow-fading loss, in dB.\\
\texttt{ue\(i\)\_rsrp} & Serving-cell RSRP approximation, in dBm.\\
\texttt{ue\(i\)\_prx} & Total serving received power, in dBm.\\
\texttt{ue\(i\)\_sinr} & Serving-link SINR, in dB.\\
\texttt{ue\(i\)\_speed} & Effective speed after pause/ramp factor, in km/h.\\
\texttt{ue\(i\)\_handover} & One if the serving index differs from the previous step.\\
\texttt{ue\(i\)\_bs\(n\)\_idx} & Index of ranked non-serving neighbor \(n\), \(1\leq n\leq6\).\\
\texttt{ue\(i\)\_bs\(n\)\_rsrp} & RSRP of that neighbor, in dBm.\\
\bottomrule
\end{longtable}

\subsection{StormSim Measurement Input Dictionary}

\begin{longtable}{p{4.3cm}p{9.6cm}}
\toprule
Column & Meaning\\
\midrule
\endfirsthead
\toprule
Column & Meaning\\
\midrule
\endhead
\texttt{Step} or \texttt{Bước} & Input sequence index.\\
\texttt{connected\_gnb} & Serving gNB identifier for the current row. A change creates a normal handover trigger.\\
\texttt{gnbX\_rsrp} & RSRP associated with gNB key \texttt{gnbX}; used by CHO or retained as measurement context.\\
\texttt{Type} & Handover type string; defaults to Xn when absent or empty.\\
\bottomrule
\end{longtable}

\subsection{Timer Event Output Dictionary}

\begin{longtable}{p{4.1cm}p{9.8cm}}
\toprule
Field & Meaning\\
\midrule
\endfirsthead
\toprule
Field & Meaning\\
\midrule
\endhead
\texttt{ts}, \texttt{start\_ts}, \texttt{stop\_ts} & Event and lifecycle timestamps.\\
\texttt{trial\_id} & Measurement-row/trial identifier.\\
\texttt{pdr} & Currently stores configured packet-loss probability despite its name.\\
\texttt{delay\_ms} & Configured delay copied into telemetry.\\
\texttt{ho\_type} & Xn or N2/NG label.\\
\texttt{timer\_name} & Timer whose lifecycle generated the event.\\
\texttt{event} & Start, stop, timeout, cancel, restart, zombie, or transport-failure marker.\\
\texttt{elapsed\_ms} & Duration measured by the timer engine.\\
\texttt{result}, \texttt{stop\_reason} & Interpreted terminal status.\\
\texttt{source\_gnb}, \texttt{target\_gnb} & Handover endpoints associated with the timer.\\
\texttt{rsrp\_at\_trigger} & Defined telemetry field; not populated in the selected logs.\\
\texttt{rlf\_detected}, beam/RLF/N2 fields & Defined for planned scenarios; zero or empty does not prove a measured zero.\\
\bottomrule
\end{longtable}

\subsection{Trial Summary Output Dictionary}

\begin{longtable}{p{4.1cm}p{9.8cm}}
\toprule
Field & Meaning\\
\midrule
\endfirsthead
\toprule
Field & Meaning\\
\midrule
\endhead
\texttt{trial\_id} & Handover attempt identifier.\\
\texttt{overall\_result} & \texttt{success} or \texttt{fail}.\\
\texttt{fail\_reason} & Scenario-level failure reason.\\
\texttt{total\_ho\_ms} & Time from monitor procedure start to observed completion/failure.\\
\texttt{source\_gnb}, \texttt{target\_gnb} & Source and requested target identifiers.\\
\texttt{pdr}, \texttt{delay\_ms}, \texttt{ho\_type} & Copied experiment-condition labels.\\
\texttt{rlf\_offset\_ms}, \texttt{rsrp\_during\_ho}, \texttt{reestab\_target\_gnb} & Planned RLF-scenario fields, generally unpopulated in the selected runs.\\
\bottomrule
\end{longtable}

\subsection{Reproducibility Checklist}

An official experiment package should retain:

\begin{enumerate}
    \item an immutable source revision;
    \item the exact command line and YAML configuration;
    \item the input measurement CSV and its schema version;
    \item all random seeds;
    \item timer values, especially T304;
    \item raw timer events and trial summaries;
    \item expected trial count and a timer-lifecycle completeness report;
    \item an explicit statement that \texttt{pdr} means loss probability or a corrected field name.
\end{enumerate}

\clearpage
